\section{The level deficits and the volume profile}\label{sec:identity}

This section is the analytic heart of the paper. We package the Saint--Venant
deficit
\[
   \dT=\dT(\Om)=\frac\pi8-T(\Om)\ \ge\ 0
\]
of a normalised domain $|\Om|=\pi$ into the level sets of its torsion
function $u$. Recall (\cref{sec:prelim}) that $u\in H^1_0(\Om)\cap
C^\infty(\Om)$ solves $-\Delta u=1$ in $\Om$, $u=0$ on $\partial\Om$, that
\[
   T(\Om):=\int_\Om u\,dx=\int_\Om|\nabla u|^2\,dx,
   \qquad m:=\max_\Om u ,
\]
the second equality holding because testing $-\Delta u=1$ against
$u\in H^1_0(\Om)$ gives $\int_\Om|\nabla u|^2=\int_\Om u$. We define two
non-negative defects of each level set; their integral over all levels turns
out to be exactly $4\pi\dT$. We then use Talenti's comparison to pin the
volume profile $\mu(v):=|\{u>v\}|$, obtaining $m\le\tfrac14$, an exact
profile identity, and the pointwise bound $-\mu'\ge4\pi$. These facts feed
directly into the inset construction of \cref{sec:inset}.

Throughout, a value $v\in(0,m)$ is called \emph{regular} if $\nabla u\ne0$ on
$\{u=v\}$; by \cref{lem:regular-values} almost every $v\in(0,m)$ is regular,
and for such $v$ the level set $\{u=v\}$ is a (closed, embedded)
$C^\infty$ curve in $\Om$ and the coarea quotient $1/|\nabla u|$ is finite and
smooth on it. We write
\[
   E_v:=\{x\in\Om:u(x)>v\},\qquad
   \mu(v):=|E_v|,\qquad
   P(v):=P(E_v)=\HH^1(\{u=v\})
\]
for the super-level set, its area, and its perimeter (the last equality
holding for regular $v$, the level set being a smooth curve).

\subsection{The flux and coarea identities}

The two structural identities below are the only place where the equation
$-\Delta u=1$ enters this section; everything afterwards is geometry of the
level sets. The first (flux) identity quantifies the source; the second
(coarea) is the differentiation formula for $\mu$. We prove both.

\begin{lemma}[Flux and coarea identities]\label{lem:flux}
For a.e.\ $v\in(0,m)$,
\begin{equation}\label{eq:flux}
   \int_{\{u=v\}}|\nabla u|\,d\HH^1=\mu(v),
\end{equation}
and for a.e.\ $v\in(0,m)$ the function $\mu$ is differentiable with
\begin{equation}\label{eq:coarea}
   -\mu'(v)=\int_{\{u=v\}}\frac{1}{|\nabla u|}\,d\HH^1 .
\end{equation}
Moreover $\mu$ is absolutely continuous and nonincreasing on $(0,m)$, with
$\mu(0^+)=|\Om|=\pi$ and $\mu(m^-)=0$.
\end{lemma}

\begin{proof}
\emph{Flux \eqref{eq:flux}.} We derive this directly from the weak
formulation, so that no boundary regularity of $\Om$ is needed. Recall
$u\in H^1_0(\Om)$ solves
\begin{equation}\label{eq:weak}
   \int_\Om\nabla u\cdot\nabla\varphi\,dx=\int_\Om\varphi\,dx
   \qquad\text{for all }\varphi\in H^1_0(\Om).
\end{equation}
Fix $0<v<m$ and, for small $h>0$ with $v+h<m$, take the Lipschitz cut-off
\[
   \psi_h(x):=\min\Bigl\{1,\ \tfrac1h\bigl(u(x)-v\bigr)_+\Bigr\}\in[0,1].
\]
Since $t\mapsto\min\{1,\tfrac1h(t-v)_+\}$ is Lipschitz and vanishes for
$t\le v$ (in particular at $t=0$), and $u\in H^1_0(\Om)$, the composition
$\psi_h$ lies in $H^1_0(\Om)$ and is an admissible test function in
\eqref{eq:weak}, with
\[
   \nabla\psi_h=\tfrac1h\,\nabla u\,\mathbf 1_{\{v<u<v+h\}}\quad\text{a.e.}
\]
Inserting $\varphi=\psi_h$ into \eqref{eq:weak},
\[
   \frac1h\int_{\{v<u<v+h\}}|\nabla u|^2\,dx
   =\int_\Om\nabla u\cdot\nabla\psi_h\,dx
   =\int_\Om\psi_h\,dx .
\]
We pass to the limit $h\downarrow0$ on each side. On the left, the coarea
formula (\cref{thm:coarea}) gives
\[
   \frac1h\int_{\{v<u<v+h\}}|\nabla u|^2\,dx
   =\frac1h\int_v^{v+h}\Bigl(\int_{\{u=t\}}|\nabla u|\,d\HH^1\Bigr)dt
   \ \xrightarrow[h\downarrow0]{}\ \int_{\{u=v\}}|\nabla u|\,d\HH^1
\]
for a.e.\ $v$, namely at every Lebesgue point of the $L^1(0,m)$ function
$t\mapsto F(t):=\int_{\{u=t\}}|\nabla u|\,d\HH^1$ (that $F\in L^1$ follows from
$\int_0^m F\,dt=\int_{\{0<u<m\}}|\nabla u|^2\,dx\le\int_\Om|\nabla u|^2<\infty$
by coarea). On the right, $0\le\psi_h\le1$ and $\psi_h\uparrow\mathbf
1_{\{u>v\}}$ pointwise as $h\downarrow0$, so by monotone (or dominated, since
$|\Om|<\infty$) convergence
\[
   \int_\Om\psi_h\,dx\ \xrightarrow[h\downarrow0]{}\ \int_\Om\mathbf
   1_{\{u>v\}}\,dx=|\{u>v\}|=\mu(v).
\]
Equating the two limits yields \eqref{eq:flux} for a.e.\ $v\in(0,m)$.

\emph{Coarea \eqref{eq:coarea} and absolute continuity.} This is the
differentiation of the distribution function, \cref{lem:dist-deriv}: for
a.e.\ $v$, $\mu$ is differentiable with $-\mu'(v)=\int_{\{u=v\}}|\nabla
u|^{-1}\,d\HH^1$. We record the absolute-continuity statement explicitly, as
it is used repeatedly. Set $\Phi(v):=\int_{\{u=v\}}|\nabla u|^{-1}\,d\HH^1\ge0$
for those $v$ where this is finite. The coarea formula (\cref{thm:coarea})
with $g=\mathbf 1_{\{u>0\}}|\nabla u|^{-1}\mathbf 1_{\{\nabla u\ne0\}}$ gives
\[
   \int_0^m\Phi(\tau)\,d\tau
   =\int_{\{0<u<m\}\cap\{\nabla u\ne0\}}1\,dx
   =|\Om|=\pi<\infty ,
\]
where the middle equality uses that $\{\nabla u=0\}\cap\{0<u<m\}$ has
Lebesgue measure zero (it is contained in the union of level sets at critical
values, a set of measure zero by Sard's theorem,
\cref{thm:sard,lem:regular-values}, intersected with $\Om$). Hence $\Phi\in
L^1(0,m)$. For $0<v_1<v_2<m$ the coarea formula applied with
$g=\mathbf 1_{\{v_1<u\le v_2\}}|\nabla u|^{-1}\mathbf 1_{\{\nabla u\ne0\}}$
yields
\[
   \mu(v_1)-\mu(v_2)=|\{v_1<u\le v_2\}|=\int_{v_1}^{v_2}\Phi(\tau)\,d\tau .
\]
Thus $\mu(v)=\int_v^m\Phi(\tau)\,d\tau$ for $v\in(0,m)$; being the tail
integral of an $L^1$ function, $\mu$ is absolutely continuous and
nonincreasing on $(0,m)$, with $\mu'=-\Phi$ a.e., which is \eqref{eq:coarea}.
Finally $\mu(0^+)=\int_0^m\Phi=|\Om|=\pi$ (equivalently $|\{u>0\}|=|\Om|$
since $u>0$ in $\Om$ by the strong minimum principle), and
$\mu(m^-)=\int_m^m\Phi=0$.
\end{proof}

\subsection{The two level deficits}

\begin{definition}[The two level deficits]\label{def:deficits}
For a regular value $v\in(0,m)$ we define the \emph{Cauchy--Schwarz} (or
H\"older/Serrin) defect and the \emph{isoperimetric} defect of the level set
$\{u=v\}$ by
\[
   D_H(v):=\mu(v)\,\bigl(-\mu'(v)\bigr)-P(v)^2 ,
   \qquad
   D_I(v):=P(v)^2-4\pi\,\mu(v).
\]
\end{definition}

The names anticipate \cref{lem:nonneg}: $D_H$ vanishes precisely when the
flux is equidistributed (constant $|\nabla u|$ on the level curve), and $D_I$
vanishes precisely when the level set is a disk.

\begin{lemma}[Non-negativity of the deficits]\label{lem:nonneg}
For a.e.\ $v\in(0,m)$ one has
\[
   D_H(v)\ \ge\ 0,\qquad D_I(v)\ \ge\ 0 .
\]
Moreover $D_I(v)=0$ if and only if $E_v$ is (equivalent to) a disk, and, for
regular $v$, $D_H(v)=0$ if and only if $|\nabla u|$ is constant
$\HH^1$-almost everywhere on $\{u=v\}$.
\end{lemma}

\begin{proof}
\emph{$D_I\ge0$.} This is the planar isoperimetric inequality
(\cref{thm:isoperimetric}) applied to the finite-perimeter set $E_v$:
$P(E_v)^2\ge4\pi|E_v|$, i.e.\ $P(v)^2\ge4\pi\mu(v)$, with equality if and
only if $E_v$ is a disk up to a Lebesgue-null modification.

\emph{$D_H\ge0$.} Fix a regular value $v$ at which \eqref{eq:flux} and
\eqref{eq:coarea} hold (a.e.\ $v$, by \cref{lem:flux,lem:regular-values}); then
$\Sigma=\{u=v\}$ is a smooth embedded curve on which $|\nabla u|>0$ and
$P(v)=\HH^1(\Sigma)<\infty$. Apply the Cauchy--Schwarz inequality
on the finite measure $\HH^1\llcorner\Sigma$ to the functions
$|\nabla u|^{1/2}$ and $|\nabla u|^{-1/2}$:
\[
   P(v)^2
   =\Bigl(\int_{\Sigma}1\,d\HH^1\Bigr)^2
   =\Bigl(\int_{\Sigma}|\nabla u|^{1/2}\cdot|\nabla u|^{-1/2}\,d\HH^1\Bigr)^2
   \le\Bigl(\int_{\Sigma}|\nabla u|\,d\HH^1\Bigr)
      \Bigl(\int_{\Sigma}\frac{1}{|\nabla u|}\,d\HH^1\Bigr).
\]
By \eqref{eq:flux} the first factor is $\mu(v)$, and by \eqref{eq:coarea} the
second factor is $-\mu'(v)$. Hence $P(v)^2\le\mu(v)(-\mu'(v))$, i.e.\
$D_H(v)\ge0$. Equality in
Cauchy--Schwarz holds if and only if $|\nabla u|^{1/2}$ and
$|\nabla u|^{-1/2}$ are proportional $\HH^1$-a.e.\ on $\Sigma$, i.e.\ if and
only if $|\nabla u|$ is $\HH^1$-a.e.\ constant on $\{u=v\}$.
\end{proof}

\subsection{The level-set deficit identity}

We now prove the central identity and, as the mandate requires, derive the
constant $4\pi$ independently of any prior computation, directly from coarea
and integration by parts (the Saint--Venant route).

\begin{theorem}[Level-set deficit identity]\label{thm:level-identity}
With the notation above,
\begin{equation}\label{eq:level-identity}
   \int_0^m\bigl(D_H(v)+D_I(v)\bigr)\,dv\ =\ 4\pi\,\dT(\Om).
\end{equation}
\end{theorem}

\begin{proof}
The proof has three steps: a pointwise telescoping, two integrations, and
the assembly with the explicit value of the constant.

\emph{Step 1: telescoping.} For a.e.\ $v\in(0,m)$ the squared-perimeter terms
in \cref{def:deficits} cancel:
\begin{equation}\label{eq:DHDI}
   D_H(v)+D_I(v)
   =\Bigl(\mu(v)(-\mu'(v))-P(v)^2\Bigr)+\Bigl(P(v)^2-4\pi\mu(v)\Bigr)
   =\mu(v)\bigl(-\mu'(v)\bigr)-4\pi\,\mu(v).
\end{equation}
The perimeter has disappeared entirely; what remains involves only $\mu$ and
$\mu'$.

\emph{Step 2: the first term, by integration by parts.} By \cref{lem:flux},
$\mu$ is absolutely continuous on $(0,m)$ with $\mu'\in L^1$, $\mu(0^+)=\pi$
and $\mu(m^-)=0$. Consequently $v\mapsto\tfrac12\mu(v)^2$ is absolutely
continuous with derivative $\mu(v)\mu'(v)\in L^1(0,m)$ (a product of a
bounded absolutely continuous function and an $L^1$ function, with the
Leibniz rule valid because $\mu$ is bounded and absolutely continuous).
Hence
\begin{equation}\label{eq:firstterm}
   \int_0^m\mu(v)\bigl(-\mu'(v)\bigr)\,dv
   =-\Bigl[\tfrac12\mu(v)^2\Bigr]_{v=0^+}^{v=m^-}
   =\tfrac12\mu(0^+)^2-\tfrac12\mu(m^-)^2
   =\tfrac12|\Om|^2=\frac{\pi^2}{2}.
\end{equation}

\emph{Step 3: the second term, by the layer-cake formula.} Since $u>0$ in
$\Om$ and $\max u=m$,
\begin{equation}\label{eq:secondterm}
   \int_0^m 4\pi\,\mu(v)\,dv
   =4\pi\int_0^\infty|\{u>v\}|\,dv
   =4\pi\int_\Om u\,dx=4\pi\,T(\Om),
\end{equation}
the middle equality being Cavalieri's (layer-cake) formula for the
nonnegative function $u$, valid since $\mu(v)=0$ for $v\ge m$.

\emph{Assembly and the constant.} Integrating \eqref{eq:DHDI} over $(0,m)$
and inserting \eqref{eq:firstterm}--\eqref{eq:secondterm},
\[
   \int_0^m\bigl(D_H+D_I\bigr)\,dv
   =\frac{\pi^2}{2}-4\pi\,T(\Om).
\]
It remains to identify $\frac{\pi^2}{2}-4\pi T(\Om)$ with $4\pi\dT$. We do
\emph{not} invoke any earlier formula for $T(B_1)$; instead we read off the
constant from Step 2. Indeed
\[
   \frac{\pi^2}{2}=\frac{|\Om|^2}{2}=4\pi\cdot\frac{|\Om|^2}{8\pi}
   =4\pi\cdot\frac{\pi^2}{8\pi}=4\pi\cdot\frac{\pi}{8},
\]
so that
\[
   \int_0^m\bigl(D_H+D_I\bigr)\,dv
   =4\pi\Bigl(\frac{\pi}{8}-T(\Om)\Bigr)=4\pi\,\dT(\Om),
\]
by the definition $\dT=\frac\pi8-T(\Om)$. This proves
\eqref{eq:level-identity}.

We pause to make the provenance of the constant transparent, since it is the
keystone of the paper. The factor $4\pi$ in $D_I$ is the planar isoperimetric
constant ($P^2\ge4\pi A$). The factor $\tfrac12|\Om|^2$ in Step 2 is purely
the integration of $\mu(-\mu')$, owing nothing to the equation beyond the
boundary values $\mu(0^+)=|\Om|$, $\mu(m^-)=0$. These two independent inputs
combine into the deficit because of the elementary algebraic identity
$\tfrac12 A^2=4\pi\cdot A^2/(8\pi)$ for any area $A$; with $A=|\Om|=\pi$ this
reads $\tfrac12|\Om|^2=4\pi\cdot\tfrac\pi8$, which is the algebra displayed
above. The number $\tfrac\pi8$ so produced is exactly the disk torsion
$T(B_1)$ (computed independently in \cref{lem:profile} below), so the
right-hand side is $4\pi\bigl(T(B_1)-T(\Om)\bigr)$. Thus
the right-hand side is $4\pi(T(B_1)-T(\Om))$, the Saint--Venant deficit
scaled by $4\pi$, with no hidden normalisation.
\end{proof}

\begin{remark}[The identity as an exact, scale-aware statement]\label{rem:identity}
Identity \eqref{eq:level-identity} expresses the (global) Saint--Venant
deficit as the total over all levels of two pointwise non-negative geometric
defects of the level sets: $D_H$, the failure of $|\nabla u|$ to be constant
on $\{u=v\}$; and $D_I$, the failure of $E_v$ to be a disk. Both vanish
identically (for all $v$) if and only if every level set is a disk on which
$|\nabla u|$ is constant, i.e.\ if and only if $\Om$ is a disk and $u=u_B$;
this is Serrin's overdetermined characterisation, and it is the rigidity
behind \eqref{eq:level-identity}. The quantitative content is the converse,
exploited from \cref{sec:inset} on: small $\dT$ forces both defects to be
small \emph{on most levels}.
\end{remark}

\subsection{Scale invariance of the identity}

For the inset of \cref{sec:inset} we will apply the identity not to $\Om$ but
to its super-level sets, which are not area-normalised. We therefore record
the unnormalised form of \eqref{eq:level-identity} and its exact
homogeneity. Let $E\subset\R^2$ be any bounded open set of finite measure
which is regular for the Dirichlet problem, with torsion function $w_E$
($-\Delta w_E=1$ in $E$, $w_E\in H^1_0(E)$), torsion
$T(E)=\int_E w_E=\int_E|\nabla w_E|^2$, maximum $m_E=\max w_E$, and level
quantities $\mu_E(v)=|\{w_E>v\}|$, $P_E(v)=P(\{w_E>v\})$,
\[
   D_H^E(v):=\mu_E(v)\bigl(-\mu_E'(v)\bigr)-P_E(v)^2,
   \qquad
   D_I^E(v):=P_E(v)^2-4\pi\mu_E(v).
\]

\begin{proposition}[Unnormalised identity and degree-$4$ homogeneity]
\label{prop:scale}
For any such $E$,
\begin{equation}\label{eq:identity-gen}
   \int_0^{m_E}\bigl(D_H^E(v)+D_I^E(v)\bigr)\,dv
   =\tfrac12|E|^2-4\pi\,T(E)
   =4\pi\bigl(T(B_E)-T(E)\bigr),
   \qquad T(B_E):=\frac{|E|^2}{8\pi},
\end{equation}
where $B_E$ is any disk with $|B_E|=|E|$. Both sides are homogeneous of
degree $4$ under dilations: writing $sE=\{sx:x\in E\}$ for $s>0$, one has
\begin{equation}\label{eq:homog}
   \int_0^{m_{sE}}\bigl(D_H^{sE}+D_I^{sE}\bigr)\,dv
   =s^4\int_0^{m_E}\bigl(D_H^{E}+D_I^{E}\bigr)\,dv,
   \qquad
   T(B_{sE})-T(sE)=s^4\bigl(T(B_E)-T(E)\bigr).
\end{equation}
Consequently the \emph{scale-invariant} deficit
$\dstar(E):=\bigl(T(B_E)-T(E)\bigr)/T(B_E)\in[0,1)$ satisfies
\begin{equation}\label{eq:scaleinv-identity}
   \dstar(E)
   =\frac{T(B_E)-T(E)}{T(B_E)}
   =\frac{\frac{1}{4\pi}\int_0^{m_E}(D_H^E+D_I^E)\,dv}{|E|^2/(8\pi)}
   =\frac{2}{|E|^2}\int_0^{m_E}\bigl(D_H^E+D_I^E\bigr)\,dv ,
\end{equation}
and $\dstar(sE)=\dstar(E)$ for all $s>0$.
\end{proposition}

\begin{proof}
\emph{Identity \eqref{eq:identity-gen}.} The proof of
\cref{lem:flux,thm:level-identity} used only $-\Delta w_E=1$, $w_E\in
H^1_0(E)$, interior smoothness, and Sard's theorem, all of which hold for $E$
verbatim (interior regularity and analyticity of the torsion function depend
only on the interior equation, not on boundary regularity;
\cref{thm:torsion-reg}). Repeating the three steps gives
\[
   \int_0^{m_E}(D_H^E+D_I^E)\,dv
   =\tfrac12\mu_E(0^+)^2-4\pi\,T(E)
   =\tfrac12|E|^2-4\pi\,T(E),
\]
since $\mu_E(0^+)=|\{w_E>0\}|=|E|$ ($w_E>0$ in $E$). The last equality in
\eqref{eq:identity-gen} is the algebraic identity
$\tfrac12|E|^2=4\pi\cdot|E|^2/(8\pi)=4\pi\,T(B_E)$.

\emph{Homogeneity \eqref{eq:homog}.} The torsion function transforms under
dilation by $w_{sE}(x)=s^2\,w_E(x/s)$: indeed $w_{sE}$ so defined vanishes on
$\partial(sE)$ and $-\Delta w_{sE}(x)=-s^2\cdot s^{-2}(\Delta w_E)(x/s)=1$,
and by uniqueness it is the torsion function of $sE$. Hence
$m_{sE}=s^2 m_E$, and
\[
   T(sE)=\int_{sE}w_{sE}\,dx=\int_{sE}s^2 w_E(x/s)\,dx
   \overset{x=sy}{=}s^2\!\int_E s^2 w_E(y)\,dy=s^4\,T(E),
\]
while $|sE|^2=s^4|E|^2$; both terms on the right of \eqref{eq:identity-gen}
scale by $s^4$, giving the second equality in \eqref{eq:homog}, and the first
follows from it via \eqref{eq:identity-gen}. (Directly: for a level $v=s^2v'$
one has $\{w_{sE}>v\}=s\{w_E>v'\}$, so $\mu_{sE}(s^2v')=s^2\mu_E(v')$,
$P_{sE}(s^2v')=s\,P_E(v')$, $\mu_{sE}'(s^2v')=\mu_E'(v')$, whence
$D_H^{sE}(s^2v')=s^2 D_H^E(v')$ and $D_I^{sE}(s^2v')=s^2 D_I^E(v')$; the
change of variables $v=s^2v'$, $dv=s^2\,dv'$ then produces the factor
$s^2\cdot s^2=s^4$.)

\emph{Scale invariance of $\dstar$.} From \eqref{eq:identity-gen},
$T(B_E)-T(E)=\tfrac{1}{4\pi}\int_0^{m_E}(D_H^E+D_I^E)\,dv$ and
$T(B_E)=|E|^2/(8\pi)$, so dividing gives
$\dstar(E)=\tfrac{2}{|E|^2}\int_0^{m_E}(D_H^E+D_I^E)\,dv$, which is
\eqref{eq:scaleinv-identity}. By \eqref{eq:homog} numerator and denominator
$T(B_E)$ both scale by $s^4$, so $\dstar(sE)=\dstar(E)$.
\end{proof}

\begin{remark}[Consistency with the normalisation]\label{rem:consistency}
Applied to $E=\Om$ with $|\Om|=\pi$, \eqref{eq:identity-gen} returns
\eqref{eq:level-identity}: $\tfrac12\pi^2-4\pi T(\Om)=4\pi\dT$, and
$T(B_\Om)=\pi^2/(8\pi)=\pi/8=T(B_1)$. The two normalisations are related by
$\dstar(\Om)=\dT/T(B_1)=\tfrac8\pi\dT$, in agreement with the global
notation.
\end{remark}

\subsection{Talenti's comparison and the volume profile}

To convert ``small on most levels'' into quantitative geometry we need the
size of the level sets, i.e.\ the profile $\mu(v)$. This is furnished by
Talenti's pointwise comparison with the model disk. Recall the unit-disk
torsion function $u_{B_1}(x)=\tfrac14(1-|x|^2)$ and its distribution function
\[
   \mu_B(v):=|\{u_{B_1}>v\}|=\pi(1-4v)_+ ,
   \qquad v\ge0,
\]
where $t_+=\max\{t,0\}$.

\begin{theorem}[Talenti comparison for the profile]\label{thm:talenti}
Let $u$ be the torsion function of $\Om$ with $|\Om|=\pi$. Then
\begin{equation}\label{eq:talenti}
   \mu(v)=|\{u>v\}|\ \le\ \mu_B(v)=\pi(1-4v)_+
   \qquad\text{for all }v\ge0 .
\end{equation}
\end{theorem}

\begin{proof}
We invoke Talenti's comparison theorem
\cite[Theorem~1]{Talenti1976} in the form proved there: if
$u\in H^1_0(\Om)$ solves $-\Delta u=f$ with $f\ge0$, and $u^\ast$ denotes the
Schwarz (radially symmetric decreasing) rearrangement of $u$, while
$\tilde u$ solves $-\Delta\tilde u=f^\ast$ in the ball $\Om^\ast$ centred at
the origin with $|\Om^\ast|=|\Om|$ and $\tilde u\in H^1_0(\Om^\ast)$, then
\[
   u^\ast(x)\ \le\ \tilde u(x)\qquad\text{for all }x\in\Om^\ast .
\]
We apply this with $f\equiv1$. Then $f^\ast\equiv1$, $\Om^\ast=B_1$ (since
$|\Om|=\pi=|B_1|$), and the comparison solution is the disk torsion function
$\tilde u=u_{B_1}$, because $u_{B_1}=\tfrac14(1-|x|^2)$ solves
$-\Delta u_{B_1}=1$ in $B_1$ with zero boundary data and is its own
rearrangement. Thus $u^\ast\le u_{B_1}$ pointwise on $B_1$.

Schwarz rearrangement is equimeasurable: $|\{u>v\}|=|\{u^\ast>v\}|$ for every
$v$. Since $u^\ast\le u_{B_1}$, we have
$\{u^\ast>v\}\subseteq\{u_{B_1}>v\}$, hence
$|\{u^\ast>v\}|\le|\{u_{B_1}>v\}|$. Combining,
\[
   \mu(v)=|\{u>v\}|=|\{u^\ast>v\}|\le|\{u_{B_1}>v\}|=\mu_B(v)
   =\pi(1-4v)_+ ,
\]
where the last equality is $\{u_{B_1}>v\}=\{|x|^2<1-4v\}$, of area
$\pi(1-4v)_+$.
\end{proof}

\begin{lemma}[Exact volume profile and the cap on the maximum]\label{lem:profile}
The Saint--Venant deficit equals the $L^1$ gap between the disk profile and
the domain profile over $(0,\tfrac14)$:
\begin{equation}\label{eq:profile}
   \int_0^{1/4}\bigl(\mu_B(v)-\mu(v)\bigr)\,dv\ =\ \dT(\Om),
\end{equation}
the integrand $\mu_B(v)-\mu(v)\ge0$ being nonnegative for every $v$. In
particular the maximum is capped,
\begin{equation}\label{eq:mcap}
   m=\max_\Om u\ \le\ \tfrac14 .
\end{equation}
\end{lemma}

\begin{proof}
\emph{The maximum cap \eqref{eq:mcap}.} Suppose, for contradiction, that
$m>\tfrac14$. Pick $v$ with $\tfrac14<v<m$. Then $\{u>v\}\ne\emptyset$ and is
open (as $u$ is continuous), so $\mu(v)=|\{u>v\}|>0$; but $v>\tfrac14$ gives
$\mu_B(v)=\pi(1-4v)_+=0$, contradicting \eqref{eq:talenti}. Hence
$m\le\tfrac14$. In particular $\mu(v)=0$ for $v\ge\tfrac14\ge m$.

\emph{The profile identity \eqref{eq:profile}.} Nonnegativity of the
integrand on all of $(0,\tfrac14)$ is exactly \eqref{eq:talenti}. By the
layer-cake formula and $m\le\tfrac14$,
\[
   T(\Om)=\int_\Om u\,dx=\int_0^\infty\mu(v)\,dv
   =\int_0^{1/4}\mu(v)\,dv ,
\]
the last step using $\mu(v)=0$ for $v\ge\tfrac14$. Likewise, by the same
formula for the disk torsion function $u_{B_1}$,
\[
   T(B_1)=\int_{B_1}u_{B_1}\,dx=\int_0^\infty\mu_B(v)\,dv
   =\int_0^{1/4}\pi(1-4v)\,dv
   =\pi\Bigl[v-2v^2\Bigr]_0^{1/4}
   =\pi\Bigl(\tfrac14-\tfrac18\Bigr)=\frac\pi8 .
\]
Subtracting,
\[
   \int_0^{1/4}\bigl(\mu_B(v)-\mu(v)\bigr)\,dv
   =T(B_1)-T(\Om)=\frac\pi8-T(\Om)=\dT(\Om).
\]
This is \eqref{eq:profile}. (The computation also reconfirms
independently that $T(B_1)=\pi/8$, hence that the deficit is well defined and
non-negative: $\dT=\int_0^{1/4}(\mu_B-\mu)\ge0$.)
\end{proof}

\begin{remark}[Two complementary decompositions of $\dT$]\label{rem:two}
We have written $\dT$ in two ways: as an integral over levels of the
\emph{shape} defects $D_H+D_I$ (\cref{thm:level-identity}, weight $4\pi$), and
as an integral over levels of the \emph{volume} gap $\mu_B-\mu$
(\cref{lem:profile}, weight $1$). The inset construction in \cref{sec:inset}
uses both: the first transfers the torsion deficit to a single level set, the
second certifies that this level set captures almost all of the area of
$\Om$.
\end{remark}

\subsection{A pointwise lower bound on \texorpdfstring{$-\mu'$}{-mu'}}

The final consequence we extract is a clean pointwise estimate, used to
quantify how fast the volume profile decreases.

\begin{lemma}[Lower bound on the profile slope]\label{lem:muprime}
For a.e.\ $v\in(0,m)$,
\begin{equation}\label{eq:muprime}
   -\mu'(v)\ \ge\ 4\pi .
\end{equation}
\end{lemma}

\begin{proof}
Fix $v\in(0,m)$ a regular value at which \eqref{eq:coarea} holds (a.e.\ $v$).
By \cref{lem:nonneg}, $D_H(v)\ge0$ and $D_I(v)\ge0$, i.e.
\[
   \mu(v)\bigl(-\mu'(v)\bigr)\ \ge\ P(v)^2\qquad\text{and}\qquad
   P(v)^2\ \ge\ 4\pi\,\mu(v).
\]
Chaining these,
\[
   \mu(v)\bigl(-\mu'(v)\bigr)\ \ge\ P(v)^2\ \ge\ 4\pi\,\mu(v).
\]
Since $v<m$, the super-level set $E_v=\{u>v\}$ is nonempty and open, so
$\mu(v)>0$; dividing by $\mu(v)>0$ gives $-\mu'(v)\ge4\pi$.
\end{proof}

\begin{remark}[Sharpness and the model]\label{rem:sharp-mu}
All four statements of this section are sharp and are simultaneously
saturated by the disk: for $u_{B_1}$ one has $\mu=\mu_B$, $-\mu_B'(v)=4\pi$ on
$(0,\tfrac14)$, $D_H\equiv D_I\equiv0$, $\dT=0$, and $m=\tfrac14$. Thus
\eqref{eq:muprime} is an equality for the model, and the level identity
\eqref{eq:level-identity} reads $0=0$. The quantitative theory of the
remaining sections measures the deviation from this rigid configuration.
\end{remark}
